\documentclass[12pt]{article}
\usepackage[utf8]{inputenc}
\usepackage[T1]{fontenc}
\usepackage{hyperref}
\title{Vao Language Resources for K-12 Curriculum Development}
\author{Compiled by Claude Sonnet 5 (claude-sonnet-5) for Andr\'as Kornai}
\date{}
\begin{document}
\maketitle

\section{Introduction}
Vao (ISO 639-3 vao) is an Austronesian language of Vanuatu, spoken on Vao island and nearby areas off the north-east coast of Malekula. It is separate from the Uripiv-Wala-Rano-Atchin language (upv), which is spoken on neighbouring islands. Searches found little material specific to Vao, and no school resources.

\section{Primary Educational Resources}
No primers, textbooks or school programmes in Vao were found in the searches.

\section{Linguistic and Reference Materials}
\begin{itemize}
\item Capell, grammar of Atchin, Pacific Linguistics D-20 (built mainly on Atchin, with comparisons to Vao; a related language, not Vao itself): \url{https://openresearch-repository.anu.edu.au/bitstream/1885/144536/1/PL-D20.pdf}
\item PARADISEC field notes series VVAO, Malekula north-east coast: \url{https://paradisec.org.au/fieldnotes/VVAO.htm}
\end{itemize}

\section{Religious and Community Resources}
\begin{itemize}
\item Joshua Project, Vanuatu people-group listing (reports a Vao Bible translation in progress; this is the only source and may be out of date): \url{https://joshuaproject.net/index.php/countries/NH?page=5}
\end{itemize}

\section{Digital and Audio Resources}
\begin{itemize}
\item PARADISEC item, collection VKS (item 69\_25): \url{https://catalog.paradisec.org.au/collections/VKS/items/69_25}
\item PARADISEC item Vao005: \url{https://catalog.paradisec.org.au/repository/ASMPI/Vao005}
\item PARADISEC item Vao007: \url{https://catalog.paradisec.org.au/repository/ASMPI/Vao007}
\end{itemize}
PARADISEC holds Vao recordings, including a 1981 narrative from Vao island and 2016 recordings from the Nalu and Torano dialect areas.

\section{Curriculum Development Context}
Archival recordings exist, but no published teaching materials were found. Mutual intelligibility with neighbouring languages was not established. A curriculum would need new materials built with the community, possibly drawing on the PARADISEC recordings subject to their access conditions.
\end{document}